Elastic Weight Consolidation (EWC) is widely used in continual learning, and standard implementations use the empirical Fisher matrix, whose justification beyond probabilistic classification remains incomplete. The empirical Fisher differs from the true Fisher, aligns with it only under specific assumptions, and is nevertheless used broadly in practice. We address this gap by deriving EWC from first principles. Starting from the continual-learning problem, we formulate old-task solution preservation as a constrained optimization problem (CLOC) and develop PL-CLOC, a population, low-order variant for deriving scalable mechanisms. Within this framework, preserving the old scalar loss value recovers exactly the empirical Fisher matrix (EF) used in EWC. This places empirical-Fisher EWC on direct optimization-theoretic footing without reliance on the true Fisher matrix. Preserving outputs up to motion along the old-task loss level set recovers the improved empirical Fisher (IEF) and leads to IEWC, a generalized and normalized sibling of EWC. EF implicitly scales local old-task sensitivity directions by the output's gradient magnitude, whereas IEF does not. We identify EWC-DR's vanishing-importance pathology as a special case of this gradient-scale effect on well-fitted data, and provide the improved empirical Fisher as the general solution. More broadly, the derivation separates the outer EWC mechanism from output geometry, enabling geometry-aware continual learning within the same regularization template. Empirically, across image classification, continual LLM/NLP, segmentation, forecasting, and diffusion, IEWC improves final average performance in two settings and, with sample weighting, improves performance across all tested settings. Together, these results motivate IEWC as a theoretical and practical successor.